% BHU PYQ -- Manually transcribed & error-corrected
% Paper: GLB-503: Sedimentary and Metamorphic Petrology
% Exam: B.Sc. (Hons.) Semester V Examination, 2022-23

\documentclass[11pt,a4paper]{article}
\usepackage[a4paper,top=2.5cm,bottom=2.5cm,left=2.8cm,right=2.8cm,headheight=14pt]{geometry}
\usepackage{amsmath,amssymb}
\usepackage[T1]{fontenc}
\usepackage[utf8]{inputenc}
\usepackage{lmodern,microtype,fancyhdr,enumitem,hyperref,lastpage,parskip}

\hypersetup{colorlinks=true,urlcolor=black,linkcolor=black,
  pdftitle={GLB-503 Sedimentary and Metamorphic Petrology -- BHU B.Sc. Sem V 2022-23}}
\pagestyle{fancy}\fancyhf{}
\renewcommand{\headrulewidth}{0.4pt}\renewcommand{\footrulewidth}{0.4pt}
\fancyhead[L]{\small   \textit{Banaras Hindu University}}
\fancyhead[R]{\small   \textit{GLB-503: Sedimentary \& Metamorphic Petrology}}
\fancyfoot[C]{\small Page  \thepage\ of \pageref{LastPage}}


\newlist{parts}{enumerate}{2}
\setlist[parts,1]{label=(\alph*),leftmargin=*,itemsep=5pt,topsep=3pt}

\newcommand{\pts}[1]{\hfill   \textit{[#1]}}


\begin{document}

\begin{center}
  {\large\bfseries BANARAS HINDU UNIVERSITY}\\[2pt]
  {\normalsize B.Sc. (Hons.) --- Semester V Examination, 2022-23}\\[6pt]
  \rule{0.8\linewidth}{0.5pt}\\[6pt]
  {\large\bfseries Paper: GLB-503 --- Sedimentary and Metamorphic Petrology}\\[4pt]
  {\normalsize Time: 3 Hours \hfill Maximum Marks: 70}\\[2pt]
  {\small Ref. No.: 052659}
\end{center}
\rule{\linewidth}{0.4pt}
\smallskip



\noindent   \textit{Instructions: Answer   \textbf{five} questions in all, selecting at least   \textbf{two} each from Sections A and B, and Section-C which is compulsory. All questions are of equal value.}
\bigskip

\bigskip


\noindent {\large\bfseries SECTION --- A}
\rule{\linewidth}{0.4pt}\smallskip




\noindent   \textbf{Question 1.} Write an essay on the process of formation of sedimentary rocks. \pts{14}

\medskip


\noindent   \textbf{Question 2.} Give the classification of sedimentary structures and, with suitable sketches, describe the inorganic primary sedimentary structures forming on the upper bedding planes. \pts{14}

\medskip


\noindent   \textbf{Question 3.} Write notes on   \textbf{any two} of the following: \pts{$2  \times 7 = 14$}

\begin{parts}
  \item Cement
  \item Cross-bedding
  \item Arkoses
\end{parts}

\bigskip


\noindent {\large\bfseries SECTION --- B}
\rule{\linewidth}{0.4pt}\smallskip




\noindent   \textbf{Question 4.} Discuss in detail the ACF and AKF diagrams and demonstrate their utilities. \pts{14}

\medskip


\noindent   \textbf{Question 5.} Discuss the phase rule and principles of metamorphic reactions. \pts{14}

\medskip


\noindent   \textbf{Question 6.} Write notes on   \textbf{any two} of the following: \pts{$2  \times 7 = 14$}

\begin{parts}
  \item Prograde metamorphism
  \item Metamorphic facies series
  \item Greenschist facies
\end{parts}

\medskip


\noindent   \textbf{Question 7.} Describe the progressive metamorphism of basic rocks in the $ \text{CaO}- \text{FeO}- \text{MgO}- \text{Al}_2 \text{O}_3- \text{SiO}_2$ system. \pts{14}

\bigskip


\noindent {\large\bfseries SECTION --- C}
\rule{\linewidth}{0.4pt}\smallskip




\noindent   \textbf{Question 8.} Fill in the blanks: \pts{$7  \times 2 = 14$}

\begin{parts}
  \item Metamorphism involves the process of \underline{\hspace{4cm}} recrystallization.
  \item Trace fossils are also known as \underline{\hspace{4cm}}.
  \item Quartzarenite sandstone indicates \underline{\hspace{4cm}} transport.
  \item In the high grades of metamorphism, \underline{\hspace{4cm}} variety of garnet is present.
  \item Imbrication is found in \underline{\hspace{4cm}}.
  \item Stromatolites are made up of \underline{\hspace{4cm}} algae.
  \item Granulite is a \underline{\hspace{4cm}} grade metamorphic rock.
\end{parts}

\vfill\rule{\linewidth}{0.4pt}

\begin{center}\small
\href{https://bhu-exam.vercel.app/}{Click here to visit our website}\\[4pt]
\href{https://me-aryan.vercel.app/}{Click here to visit our contributor}\\[4pt]
\href{https://quashan-me.vercel.app/}{Click here to visit our contributor}
\end{center}
\end{document}
